% =====================================================================
\chapter{Funciones trigonométricas, inversas y ecuaciones}
\label{tema:4}
% =====================================================================

Este tema estudia el seno, el coseno y la tangente como funciones de una variable real: su periodicidad, sus gráficas y las transformaciones que producen las sinusoides generales. Después restringe sus dominios para obtener funciones inversas y, con ellas, describe todas las soluciones de las ecuaciones trigonométricas. Se sigue el capítulo~8 de~\cite{osat} (§8.1 a §8.3) y la §9.5 del mismo libro; la conexión con la codificación de la posición procede de Vaswani \emph{et al.}~\cite[§3.5]{vaswani}.

\section{Funciones periódicas}

\begin{definicion}[Función periódica]
Una función $f\colon D\to\R$ es \emph{periódica} de periodo $T>0$ si, para todo $x\in D$, se cumple $x+T\in D$, $x-T\in D$ y $f(x+T)=f(x)$. Si existe un periodo mínimo, se llama \emph{periodo fundamental}.
\end{definicion}

\begin{proposicion}\label{prop:periodos}
(i) El seno y el coseno tienen periodo fundamental $2\pi$; la tangente, $\pi$. (ii) Si $f$ tiene periodo fundamental $T$ y $\omega>0$, la función $g(x)=f(\omega x)$ tiene periodo fundamental $T/\omega$.
\end{proposicion}

\section{Gráficas y sinusoides}

La \cref{fig:graficas} muestra las gráficas del seno y del coseno. Ambas oscilan entre $-1$ y $1$, y la del coseno es la del seno desplazada $\pi/2$ hacia la izquierda, porque $\cos x=\sin(x+\frac{\pi}{2})$ (\cref{teo:simetrias}(vi) con la paridad). La tangente tiene periodo $\pi$, ceros en $k\pi$ y asíntotas verticales en $\frac{\pi}{2}+k\pi$, donde se anula el coseno.

\addtocounter{figure}{1}%

\begin{definicion}[Sinusoide]
Una \emph{sinusoide} es una función $y=A\sin\bigl(\omega(x-x_0)\bigr)+D$ con $A\ne0$ y $\omega>0$. Se llaman \emph{amplitud} a $\abs{A}$, \emph{periodo} a $2\pi/\omega$, \emph{frecuencia angular} a $\omega$, \emph{desfase} (desplazamiento horizontal) a $x_0$ y \emph{desplazamiento vertical} a $D$.
\end{definicion}

Por la \cref{prop:periodos}(ii), el periodo de la sinusoide es efectivamente $2\pi/\omega$. Al escribir $y=A\sin(\omega x-c)+D$, el desfase es $x_0=c/\omega$, no $c$.

\begin{proposicion}[Suma de un seno y un coseno]\label{prop:armonico}
Sean $a,b\in\R$ no ambos nulos y $R=\sqrt{a^{2}+b^{2}}$. Existe $\varphi\in(-\pi,\pi]$ con $\cos\varphi=a/R$ y $\sin\varphi=b/R$, y para todo $x$
\[
a\sin x+b\cos x=R\sin(x+\varphi).
\]
En particular, el máximo de $a\sin x+b\cos x$ es $R$ y el mínimo, $-R$.
\end{proposicion}

\section{Funciones trigonométricas inversas}

Una función periódica no es inyectiva, luego no tiene inversa. Se restringe a un intervalo donde sea biyectiva sobre su imagen~\cite[§8.3]{osat}.

\begin{definicion}[Inversas principales]\label{def:inversas}
\begin{enumerate}
  \item $\sin\colon[-\frac{\pi}{2},\frac{\pi}{2}]\to[-1,1]$ es biyectiva; su inversa es $\arcsin\colon[-1,1]\to[-\frac{\pi}{2},\frac{\pi}{2}]$.
  \item $\cos\colon[0,\pi]\to[-1,1]$ es biyectiva; su inversa es $\arccos\colon[-1,1]\to[0,\pi]$.
  \item $\tan\colon(-\frac{\pi}{2},\frac{\pi}{2})\to\R$ es biyectiva; su inversa es $\arctan\colon\R\to(-\frac{\pi}{2},\frac{\pi}{2})$.
\end{enumerate}
\end{definicion}

La inyectividad de estas restricciones se demuestra en el \cref{eje:4-creciente} (el seno es estrictamente creciente en $[-\frac{\pi}{2},\frac{\pi}{2}]$) y de forma análoga para las otras dos. La sobreyectividad es geométrica: cada recta horizontal $y=t$ con $\abs{t}\le1$ corta la semicircunferencia derecha en un punto, y cada recta vertical $x=t$ corta la superior en un punto; para la tangente, cada recta por el origen no vertical corta la semicircunferencia derecha abierta.

\begin{proposicion}[Propiedades de las inversas]\label{prop:inversas}
\begin{enumerate}
  \item $\sin(\arcsin x)=x$ para $x\in[-1,1]$, y $\arcsin(\sin\theta)=\theta$ si y solo si $\theta\in[-\frac{\pi}{2},\frac{\pi}{2}]$.
  \item $\cos(\arcsin x)=\sqrt{1-x^{2}}$ y $\sin(\arccos x)=\sqrt{1-x^{2}}$ para $x\in[-1,1]$.
  \item $\arcsin x+\arccos x=\frac{\pi}{2}$ para $x\in[-1,1]$.
  \item $\arcsin$ y $\arctan$ son impares, y $\arccos(-x)=\pi-\arccos x$.
\end{enumerate}
\end{proposicion}

\addtocounter{definicion}{1}%

\section{Ecuaciones trigonométricas}

\begin{teorema}[Soluciones generales]\label{teo:ecuaciones}
Sea $a\in\R$.
\begin{enumerate}
  \item Si $\abs{a}\le1$ y $\alpha=\arcsin a$, entonces $\sin x=a$ si y solo si $x=\alpha+2k\pi$ o $x=\pi-\alpha+2k\pi$ con $k\in\Z$.
  \item Si $\abs{a}\le1$ y $\beta=\arccos a$, entonces $\cos x=a$ si y solo si $x=\pm\beta+2k\pi$ con $k\in\Z$.
  \item Para todo $a$, $\tan x=a$ si y solo si $x=\arctan a+k\pi$ con $k\in\Z$.
\end{enumerate}
Si $\abs{a}>1$, las ecuaciones $\sin x=a$ y $\cos x=a$ no tienen solución.
\end{teorema}

\addtocounter{definicion}{1}%

\addtocounter{definicion}{1}%

\addtocounter{definicion}{1}%

